\documentclass{zhvt-classic}

% 使用本地源流明体（GenRyuMin2 TW）按文件路径加载，不依赖系统字体名
% 版框（内框／界栏／书口侧栏／鱼尾）改用与版心印章同一印泥棕色：
%   印章图片 yshs.png 实色 = RGB(135,68,52) = HTML 874434
\usepackage{xcolor}
\usepackage{tikz}
% sealbrown 双注册：xcolor 供 TikZ/pgf，L3 供 cls 版框
\definecolor{sealbrown}{HTML}{874434}
\ExplSyntaxOn
\color_set:nnn { sealbrown } { HTML } { 874434 }
\ExplSyntaxOff

\zhvtset{
  grid color       = sealbrown,
  main font        = GenRyuMin2TW-SB,
  main font path   = C:/Users/Administrator/texmf/fonts/opentype/public/GenRyuMin2TW/,
jiazhu font      = GenRyuMin2TW-M,
  jiazhu font path = C:/Users/Administrator/texmf/fonts/opentype/public/GenRyuMin2TW/,
  % 页码框容纳 4 字：正常页「九十七」+前置空 1 格=4 行恰好；
  % 100+ 页「一百零一」共 5 行超出，末字「一/二/三」被框裁掉（2026-10-04 实测），故放宽到 5
  shukou page number chars = 5,
  % 书口书名框高从默认 6 字减到 4 字，把省出的 2 字高拨给 chapter 框，
  % 使 6 字书名+卷名放进 9 字高的 chapter 框后不再压线（2026-10-05）
  shukou title chars = 4,
}

\title{史记}[范雎蔡泽列传]
\author{汉·司马迁}
\maker{闲叟漫记·重製}


% ===== 索隐述赞样式（2026-11-08 四次定稿 · 愚大人亲手改定参数）=====
% ⚠ 竖排前提：自定义字体家族必须带 RawFeature={vertical:+vert}；
%   simhei.ttf 加 +vhal 会整体加载失败静默回退，只写 +vert；simfang 同理。
\newCJKfontfamily{\shuzanfont}[Path=C:/Windows/Fonts/,RawFeature={vertical:+vert}]{simhei.ttf}
\newCJKfontfamily{\shuzanbody}[Path=C:/Windows/Fonts/,RawFeature={vertical:+vert}]{simfang.ttf}
% 题头：152x30pt 圆角 6pt 棕底白字；baseline (0,2) 使整框落入界行格子不压版框线；
% 白字节点 x 不补偿、y +8.2pt（像素级墨迹实测 dx/dy≈0，L/R 留白约 20pt）。
\newcommand{\shuzanhead}[1]{%
  \begin{tikzpicture}[baseline={(0,2)}]
    \fill[sealbrown, rounded corners=6pt]
      (-76pt,-15pt) rectangle (76pt,15pt);
    \node[text=white, inner sep=0pt, anchor=center] at (0pt,8.2pt)
      {\shuzanfont\fontsize{23}{23}\selectfont #1};
  \end{tikzpicture}%
}
% ★ 修正书口卷名位置：chapter 框起始位置改回 cls 默认值，避免全局影响。
%   只在每个卷首页通过 \juan 命令里的空命令微调该页位置（安全，不会跨页污染）。
\ExplSyntaxOn
\AtBeginDocument{
  \dim_set:Nn \l__zhvt_shukou_chapter_title_start_dim
    { \l__zhvt_shukou_title_height_dim + \l__zhvt_font_size_dim * 2 }
  \dim_set:Nn \l__zhvt_shukou_page_number_start_dim
    { \l__zhvt_shukou_chapter_title_start_dim + \l__zhvt_shukou_chapter_title_height_dim }
}
\ExplSyntaxOff

\begin{document}

% ★★ 版心（书口）排布覆写：把书名放到鱼尾之下（2026-10-04）
%
%   cls 默认版心自上而下：书名（title 静态框，在鱼尾之上）→ 鱼尾 → 篇名（chapter 动态框）→ 页码
%   本模板改为：          鱼尾（位置不动）→ 书名 → 篇名 → 页码分割线 → 页码
%
%   手法：把 cls 放在 title 静态框里的书名，改放进 chapter 动态框（鱼尾之下）；
%         title 静态框留空。书名取 \l__zhvt_title_str，即 \title{...} 的第一参数，
%         所以换书时只改 \title，不必在此重复写书名。
%
%   ⚠ 三条必须遵守的写法（2026-10-04 实测）：
%     ① 汉字必须用 \hbox{...} 包住。flowfram 会把框内容装进一个宽 14.4mm（= 一个 line_sep
%        = 书口栏宽）的纵向 minipage 并按该宽度折行，而整页靠 PDF /Rotate 90 显示、
%        框内文字本身不变向 —— 不包 \hbox 时「围炉夜话」4 字会被逐字折成 4 行，
%        在版面上就成了横排 4 字并横向溢出版框。
%     ② 竖排下**一个 \hbox 就是一整列**。书名与篇名要同处一列，必须放进同一个 \hbox；
%        分成两个 \hbox 时，后一个会被当成新的一列、甩到隔壁栏去
%        （实测「序」跑到书口栏外侧 14.4mm 的正文列上）。中间的 \quad 在框内是横向空白，
%        显示出来正好是书名与篇名之间的纵向间隔。
%     ③ ★ 内容盒要写成 **与框等宽的定宽盒** `\hbox to \l__zhvt_shukou_chapter_title_height_dim{…\hfil}`。
%        实测：chapter 框是把内容盒沿纵向**居中**摆放的，内容宽度一变，书名就跟着漂：
%        自然宽度的 \hbox（书名 38.4mm）落在 x=157.2mm，加篇名后（57.6mm）落在 x=152.4mm，
%        目录页（无篇名）与正文页的书名就差半格。做成等宽盒后，无论有没有篇名，
%        内容固定从框顶起排，实测两页均落在 x=142.8mm，不再漂。
\ExplSyntaxOn
% 把书口 chapter 框（书名+篇名）起始位置下移 1 字高，避免 4 字卷名顶到鱼尾
% ⚠ cls 的 \__zhvt_calculate_dimensions: 在 \AtBeginDocument 里执行，会重算该 dim；
%   必须把覆写挂到 begin document 之后，否则会被二次覆盖（2026-10-05 实测）。
% (旧方案已移到导言区，直接修改 chapter 框起始位置)
% 卷标题命令 #1=书口篇名(鱼尾) #2=正文卷名（含卷序，卷名内空格用 ~）
%
% ★★ 目录页两项修订（2026-10-05，愚大人指出）★★
%   ① 页码起点统一对齐：
%      原状 —— 目录项写成 \par\hspace{4em}卷名\hspace{4em}页码，页码起点 = 卷名实宽 + 4em。
%      卷名实宽随字数变：「卷一~梁惠王上」6 字 +~、「卷七~离娄上」5 字 +~，
%      差 1 个字位 → 卷七~卷十 4 条的页码整体比其余 10 条高 27.21pt（一个字），看着上下不齐。
%      改法 —— 卷名补全角空格（宽 = 1 字 = 27.215pt）补足到 **7 字符位**再写进目录，
%      14 条页码起点遂全部落在同一条竖线上（起点对齐；末字随页码字数自然参差）。
%      ⚠ 只对目录里的副本 \l_tmpc_tl 补，正文卷名仍用 \l_tmpa_tl 原样 —— 正文像素不受影响。
%   ② 卷名整体上移半字位：前导 \hspace 由 4em 改 3.5em。
%      实测原状目录块在版框内偏下（天头留白 139.6pt / 地脚留白 110.7pt），
%      上移 0.5em = 13.6pt 后上下基本均分，视觉居中。
\NewDocumentCommand { \juan } { mm } {
	\clearpage
	\tl_set:Nn \l__zhvt_chapter_title_tl { #1 }
	\tl_set:Nn \l_tmpa_tl { #2 }
	\tl_set_eq:NN \l_tmpc_tl \l_tmpa_tl
	\int_step_inline:nn { 7 - \str_count:n { #2 } }
	  { \tl_put_right:Nn \l_tmpc_tl { 　 } }
	\tl_set:Nx \l_tmpb_tl { \zhnum { page } }
	\tl_put_left:Nn \l_tmpb_tl { \hspace { 4em } }
	\tl_put_left:Nn \l_tmpb_tl { \l_tmpc_tl }
	\tl_put_left:Nn \l_tmpb_tl { \par \hspace { 3.5em } }
	\addtocontents { toc } { \l_tmpb_tl }
	\skip_horizontal:n { \l__zhvt_font_size_dim * 4 }
	\l_tmpa_tl
	\par
	\thispagestyle { empty }
}
\cs_set_protected:Npn \__zhvt_set_shukou_content: {
	% ★ 鱼尾上方留白（title 静态框）嵌入闲叟钤印：
	%   整页 /Rotate 90，故此处把图用 \rotatebox{90} 转正为竖章；
	%   尺寸卡在版心条幅内（书内约 7mm 宽 × 33.8mm 高，留余量不压鱼尾）。
	\setstaticcontents * { title } {
		\vbox_to_ht:nn { \l__zhvt_baseline_skip_dim } {
			\vskip 14.2pt
			\hbox_to_wd:nn { \l__zhvt_shukou_title_height_dim } {
				\hfil
				\rotatebox [ origin = c ] { 90 } { \includegraphics [ width = 7mm , height = 33.8mm ] { yshs.png } }
				\hfil
			}
		}
	}
% ★ 版心（书口）统一改法（2026-11-16 愚大人定）：18pt＋框内居中＋书口文字同版框色
%   ① 原 27.2pt 时 7 字长篇名（如「廉颇蔺相如列传」）末字溢出 chapter 框、压住页码；
%      改 \fontsize{18pt}{18pt}，前后 \hfil 令字串在框内居中（不随篇名字数漂移）。
%   ② \vbox{\vskip 8.2pt} 定量下移，抵消缩字号后栏宽方向 8.2pt 上偏（像素实测 y 中心 543.3 = 原版）。
%      ⚠ \raisebox / \vspace* 在 flowfram 框内无效；换字号须同步换补偿：16pt→10.1、19pt→7.4。
%   ③ chapter 与 page~number 均置 \color{sealbrown}（HTML 874434），与版框／界行同色；正文仍黑。
%   ⚠ 页码字位、钤印、鱼尾、界行、正文一律未动；五部合订本同此口径。
	\setdynamiccontents * { chapter }
	{ \color{sealbrown} \vbox { \vskip 8.2pt \hbox to \l__zhvt_shukou_chapter_title_height_dim
		{ \hfil \fontsize{18pt}{18pt}\selectfont \l__zhvt_title_str \quad \l__zhvt_chapter_title_tl \hfil } } }
	\setdynamiccontents * { page~number } { \color{sealbrown} \quad \zhnum { page } }
}
\__zhvt_set_shukou_content:
\ExplSyntaxOff

\maketitle

% 单篇，无多卷目录，不渲染目录页

\juan{范雎蔡泽列传}{卷七十九~范雎蔡泽列传第十九}
\hspace*{2em}范雎者。魏人也。字叔。游说诸侯。欲事魏王。家贫无以自资。乃先事魏中大夫须贾。\\
\hspace*{2em}须贾为魏昭王使於齐。范雎从。留数月。未得报。齐襄王闻雎辩口。乃使人赐雎金十斤及牛酒。雎辞谢不敢受。须贾知之。大怒。以为雎持魏国阴事告齐。故得此馈。令雎受其牛酒。还其金。既归。心怒雎。以告魏相。魏相。魏之诸公子。曰魏齐。魏齐大怒。使舍人笞击雎。折胁摺齿。雎详死。即卷以箦。置厕中。宾客饮者醉。更溺雎。故僇辱以惩後。令无妄言者。雎从箦中谓守者曰。﹁公能出我。我必厚谢公。﹂守者乃请出弃箦中死人。魏齐醉。曰。﹁可矣。﹂范雎得出。後魏齐悔。复召求之。魏人郑安平闻之。乃遂操范雎亡。伏匿。更名姓曰张禄。\\
\hspace*{2em}当此时。秦昭王使谒者王稽於魏。郑安平诈为卒。侍王稽。王稽问。﹁魏有贤人可与俱西游者乎。﹂郑安平曰。﹁臣里中有张禄先生。欲见君。言天下事。其人有仇。不敢昼见。﹂王稽曰。﹁夜与俱来。﹂郑安平夜与张禄见王稽。语未究。王稽知范雎贤。谓曰。﹁先生待我於三亭之南。﹂与私约而去。\\
\hspace*{2em}王稽辞魏去。过载范雎入秦。至湖。望见车骑从西来。范雎曰。﹁彼来者为谁。﹂王稽曰。﹁秦相穰侯东行县邑。﹂范雎曰。﹁吾闻穰侯专秦权。恶内诸侯客。此恐辱我。我宁且匿车中。﹂有顷。穰侯果至。劳王稽。因立车而语曰。﹁关东有何变。﹂曰。﹁无有。﹂又谓王稽曰。﹁谒君得无与诸侯客子俱来乎。无益。徒乱人国耳。﹂王稽曰。﹁不敢。﹂即别去。范雎曰。﹁吾闻穰侯智士也。其见事迟。乡者疑车中有人。忘索之。﹂於是范雎下车走。曰。﹁此必悔之。﹂行十馀里。果使骑还索车中。无客。乃已。王稽遂与范雎入咸阳。\\
\hspace*{2em}已报使。因言曰。﹁魏有张禄先生。天下辩士也。曰﹃秦王之国危於累卵。得臣则安。然不可以书传也﹄。臣故载来。﹂秦王弗信。使舍食草具。待命岁馀。\\
\hspace*{2em}当是时。昭王已立三十六年。南拔楚之鄢郢。楚怀王幽死於秦。秦东破齐。湣王尝称帝。後去之。数困三晋。厌天下辩士。无所信。\\
\hspace*{2em}穰侯。华阳君。昭王母宣太后之弟也。而泾阳君。高陵君皆昭王同母弟也。穰侯相。三人者更将。有封邑。以太后故。私家富重於王室。及穰侯为秦将。且欲越韩。魏而伐齐纲寿。欲以广其陶封。范雎乃上书曰。\\
\hspace*{2em}臣闻明主立政。有功者不得不赏。有能者不得不官。劳大者其禄厚。功多者其爵尊。能治众者其官大。故无能者不敢当职焉。有能者亦不得蔽隐。使以臣之言为可。原行而益利其道。以臣之言为不可。久留臣无为也。语曰。﹁庸主赏所爱而罚所恶。明主则不然。赏必加於有功。而刑必断於有罪。﹂今臣之胸不足以当椹质。而要不足以待斧钺。岂敢以疑事尝试於王哉。虽以臣为贱人而轻辱。独不重任臣者之无反复於王邪。\\
\hspace*{2em}且臣闻周有砥砨。宋有结绿。梁有县藜。楚有和朴。此四宝者。土之所生。良工之所失也。而为天下名器。然则圣王之所弃者。独不足以厚国家乎。\\
\hspace*{2em}臣闻善厚家者取之於国。善厚国者取之於诸侯。天下有明主则诸侯不得擅厚者。何也。为其割荣也。良医知病人之死生。而圣主明於成败之事。利则行之。害则舍之。疑则少尝之。虽舜禹复生。弗能改已。语之至者。臣不敢载之於书。其浅者又不足听也。意者臣愚而不概於王心邪。亡其言臣者贱而不可用乎。自非然者。臣原得少赐游观之间。望见颜色。一语无效。请伏斧质。\\
\hspace*{2em}於是秦昭王大说。乃谢王稽。使以传车召范雎。\\
\hspace*{2em}於是范雎乃得见於离宫。详为不知永巷而入其中。王来而宦者怒。逐之。曰。﹁王至。﹂范雎缪为曰。﹁秦安得王。秦独有太后。穰侯耳。﹂欲以感怒昭王。昭王至。闻其与宦者争言。遂延迎。谢曰。﹁寡人宜以身受命久矣。会义渠之事急。寡人旦暮自请太后。今义渠之事已。寡人乃得受命。窃闵然不敏。敬执宾主之礼。﹂范雎辞让。是日观范雎之见者。群臣莫不洒然变色易容者。\\
\hspace*{2em}秦王屏左右。宫中虚无人。秦王跽而请曰。﹁先生何以幸教寡人。﹂范雎曰。﹁唯唯。﹂有间。秦王复跽而请曰。﹁先生何以幸教寡人。﹂范雎曰。﹁唯唯。﹂若是者三。秦王跽曰。﹁先生卒不幸教寡人邪。﹂范雎曰。﹁非敢然也。臣闻昔者吕尚之遇文王也。身为渔父而钓於渭滨耳。若是者。交疏也。已说而立为太师。载与俱归者。其言深也。故文王遂收功於吕尚而卒王天下。乡使文王疏吕尚而不与深言。是周无天子之德。而文武无与成其王业也。今臣羁旅之臣也。交疏於王。而所原陈者皆匡君之事。处人骨肉之间。原效愚忠而未知王之心也。此所以王三问而不敢对者也。臣非有畏而不敢言也。臣知今日言之於前而明日伏诛於後。然臣不敢避也。大王信行臣之言。死不足以为臣患。亡不足以为臣忧。漆身为厉被发为狂不足以为臣耻。且以五帝之圣焉而死。三王之仁焉而死。五伯之贤焉而死。乌获。任鄙之力焉而死。成荆。孟贲。王庆忌。夏育之勇焉而死。死者。人之所必不免也。处必然之势。可以少有补於秦。此臣之所大原也。臣又何患哉。伍子胥橐载而出昭关。夜行昼伏。至於陵水。无以餬其口。行蒲伏。稽首肉袒。鼓腹吹篪。乞食於吴市。卒兴吴国。阖闾为伯。使臣得尽谋如伍子胥。加之以幽囚。终身不复见。是臣之说行也。臣又何忧。箕子。接舆漆身为厉。被发为狂。无益於主。假使臣得同行於箕子。可以有补於所贤之主。是臣之大荣也。臣有何耻。臣之所恐者。独恐臣死之後。天下见臣之尽忠而身死。因以是杜口裹足。莫肯乡秦耳。足下上畏太后之严。下惑於奸臣之态。居深宫之中。不离阿保之手。终身迷惑。无与昭奸。大者宗庙灭覆。小者身以孤危。此臣之所恐耳。若夫穷辱之事。死亡之患。臣不敢畏也。臣死而秦治。是臣死贤於生。﹂秦王跽曰。﹁先生是何言也。夫秦国辟远。寡人愚不肖。先生乃幸辱至於此。是天以寡人慁先生而存先王之宗庙也。寡人得受命於先生。是天所以幸先王。而不弃其孤也。先生柰何而言若是。事无小大。上及太后。下至大臣。原先生悉以教寡人。无疑寡人也。﹂范雎拜。秦王亦再拜\\
\hspace*{2em}范雎曰。﹁大王之国。四塞以为固。北有甘泉。谷口。南带泾。渭。右陇。蜀。左关。阪。奋击百万。战车千乘。利则出攻。不利则入守。此王者之地也。民怯於私斗而勇於公战。此王者之民也。王并此二者而有之。夫以秦卒之勇。车骑之众。以治诸侯。譬若施韩卢而搏蹇兔也。霸王之业可致也。而群臣莫当其位。至今闭关十五年。不敢窥兵於山东者。是穰侯为秦谋不忠。而大王之计有所失也。﹂秦王跽曰。﹁寡人原闻失计。﹂\\
\hspace*{2em}然左右多窃听者。范雎恐。未敢言内。先言外事。以观秦王之俯仰。因进曰。﹁夫穰侯越韩。魏而攻齐纲寿。非计也。少出师则不足以伤齐。多出师则害於秦。臣意王之计。欲少出师而悉韩。魏之兵也。则不义矣。今见与国之不亲也。越人之国而攻。可乎。其於计疏矣。且昔齐湣王南攻楚。破军杀将。再辟地千里。而齐尺寸之地无得焉者。岂不欲得地哉。形势不能有也。诸侯见齐之罢弊。君臣之不和也。兴兵而伐齐。大破之。士辱兵顿。皆咎其王。曰。﹃谁为此计者乎。﹄王曰。﹃文子为之。﹄大臣作乱。文子出走。攻齐所以大破者。以其伐楚而肥韩。魏也。此所谓借贼兵而赍盗粮者也。王不如远交而近攻。得寸则王之寸也。得尺亦王之尺也。今释此而远攻。不亦缪乎。且昔者中山之国地方五百里。赵独吞之。功成名立而利附焉。天下莫之能害也。今夫韩。魏。中国之处而天下之枢也。王其欲霸。必亲中国以为天下枢。以威楚。赵。楚彊则附赵。赵彊则附楚。楚。赵皆附。齐必惧矣。齐惧。必卑辞重币以事秦。齐附而韩。魏因可虏也。﹂昭王曰。﹁吾欲亲魏久矣。而魏多变之国也。寡人不能亲。请问亲魏柰何。﹂对曰。﹁王卑词重币以事之。不可。则割地而赂之。不可。因举兵而伐之。﹂王曰。﹁寡人敬闻命矣。﹂乃拜范雎为客卿。谋兵事。卒听范雎谋。使五大夫绾伐魏。拔怀。後二岁。拔邢丘。\\
\hspace*{2em}客卿范雎复说昭王曰。﹁秦韩之地形。相错如绣。秦之有韩也。譬如木之有蠹也。人之有心腹之病也。天下无变则已。天下有变。其为秦患者孰大於韩乎。王不如收韩。﹂昭王曰。﹁吾固欲收韩。韩不听。为之柰何。﹂对曰。﹁韩安得无听乎。王下兵而攻荥阳。则巩。成皋之道不通。北断太行之道。则上党之师不下。王一兴兵而攻荥阳。则其国断而为三。夫韩见必亡。安得不听乎。若韩听。而霸事因可虑矣。﹂王曰。﹁善。﹂且欲发使於韩。\\
\hspace*{2em}范雎日益亲。复说用数年矣。因请间说曰。﹁臣居山东时。闻齐之有田文。不闻其有王也。闻秦之有太后。穰侯。华阳。高陵。泾阳。不闻其有王也。夫擅国之谓王。能利害之谓王。制杀生之威之谓王。今太后擅行不顾。穰侯出使不报。华阳。泾阳等击断无讳。高陵进退不请。四贵备而国不危者。未之有也。为此四贵者下。乃所谓无王也。然则权安得不倾。令安得从王出乎。臣闻善治国者。乃内固其威而外重其权。穰侯使者操王之重。决制於诸侯。剖符於天下。政適伐国。莫敢不听。战胜攻取则利归於陶。国弊御於诸侯。战败则结怨於百姓。而祸归於社稷。诗曰﹃木实繁者披其枝。披其枝者伤其心。大其都者危其国。尊其臣者卑其主﹄。崔杼。淖齿管齐。射王股。擢王筋。县之於庙梁。宿昔而死。李兑管赵。囚主父於沙丘。百日而饿死。今臣闻秦太后。穰侯用事。高陵。华阳。泾阳佐之。卒无秦王。此亦淖齿。李兑之类也。且夫三代所以亡国者。君专授政。纵酒驰骋弋猎。不听政事。其所授者。妒贤嫉能。御下蔽上。以成其私。不为主计。而主不觉悟。故失其国。今自有秩以上至诸大吏。下及王左右。无非相国之人者。见王独立於朝。臣窃为王恐。万世之後。有秦国者非王子孙也。﹂昭王闻之大惧。曰。﹁善。﹂於是废太后。逐穰侯。高陵。华阳。泾阳君於关外。秦王乃拜范雎为相。收穰侯之印。使归陶。因使县官给车牛以徙。千乘有馀。到关。关阅其宝器。宝器珍怪多於王室。\\
\hspace*{2em}秦封范雎以应。号为应侯。当是时。秦昭王四十一年也。\\
\hspace*{2em}范雎既相秦。秦号曰张禄。而魏不知。以为范雎已死久矣。魏闻秦且东伐韩。魏。魏使须贾於秦。范雎闻之。为微行。敝衣间步之邸。见须贾。须贾见之而惊曰。﹁范叔固无恙乎。﹂范雎曰。﹁然。﹂须贾笑曰。﹁范叔有说於秦邪。﹂曰。﹁不也。雎前日得过於魏相。故亡逃至此。安敢说乎。﹂须贾曰。﹁今叔何事。﹂范雎曰﹁臣为人庸赁。﹂须贾意哀之。留与坐饮食。曰。﹁范叔一寒如此哉。﹂乃取其一綈袍以赐之。须贾因问曰。﹁秦相张君。公知之乎。吾闻幸於王。天下之事皆决於相君。今吾事之去留在张君。孺子岂有客习於相君者哉。﹂范雎曰。﹁主人翁习知之。唯雎亦得谒。雎请为见君於张君。﹂须贾曰。﹁吾马病。车轴折。非大车驷马。吾固不出。﹂范雎曰。﹁原为君借大车驷马於主人翁。﹂\\
\hspace*{2em}范雎归取大车驷马。为须贾御之。入秦相府。府中望见。有识者皆避匿。须贾怪之。至相舍门。谓须贾曰。﹁待我。我为君先入通於相君。﹂须贾待门下。持车良久。问门下曰。﹁范叔不出。何也。﹂门下曰。﹁无范叔。﹂须贾曰。﹁乡者与我载而入者。﹂门下曰。﹁乃吾相张君也。﹂须贾大惊。自知见卖。乃肉袒行。因门下人谢罪。於是范雎盛帷帐。待者甚众。见之。须贾顿首言死罪。曰。﹁贾不意君能自致於青云之上。贾不敢复读天下之书。不敢复与天下之事。贾有汤镬之罪。请自屏於胡貉之地。唯君死生之。﹂范雎曰。﹁汝罪有几。﹂曰。﹁擢贾之发以续贾之罪。尚未足。﹂范雎曰。﹁汝罪有三耳。昔者楚昭王时而申包胥为楚卻吴军。楚王封之以荆五千户。包胥辞不受。为丘墓之寄於荆也。今雎之先人丘墓亦在魏。公前以雎为有外心於齐而恶雎於魏齐。公之罪一也。当魏齐辱我於厕中。公不止。罪二也。更醉而溺我。公其何忍乎。罪三矣。然公之所以得无死者。以綈袍恋恋。有故人之意。故释公。﹂乃谢罢。入言之昭王。罢归须贾。\\
\hspace*{2em}须贾辞於范雎。范雎大供具。尽请诸侯使。与坐堂上。食饮甚设。而坐须贾於堂下。置莝豆其前。令两黥徒夹而马食之。数曰。﹁为我告魏王。急持魏齐头来。不然者。我且屠大梁。﹂须贾归。以告魏齐。魏齐恐。亡走赵。匿平原君所。\\
\hspace*{2em}范雎既相。王稽谓范雎曰。﹁事有不可知者三。有不柰何者亦三。宫车一日晏驾。是事之不可知者一也。君卒然捐馆舍。是事之不可知者二也。使臣卒然填沟壑。是事之不可知者三也。宫车一日晏驾。君虽恨於臣。无可柰何。君卒然捐馆舍。君虽恨於臣。亦无可柰何。使臣卒然填沟壑。君虽恨於臣。亦无可柰何。﹂范雎不怿。乃入言於王曰。﹁非王稽之忠。莫能内臣於函谷关。非大王之贤圣。莫能贵臣。今臣官至於相。爵在列侯。王稽之官尚止於谒者。非其内臣之意也。﹂昭王召王稽。拜为河东守。三岁不上计。又任郑安平。昭王以为将军。范雎於是散家财物。尽以报所尝困戹者。一饭之德必偿。睚眦之怨必报。\\
\hspace*{2em}范雎相秦二年。秦昭王之四十二年。东伐韩少曲。高平。拔之。\\
\hspace*{2em}秦昭王闻魏齐在平原君所。欲为范雎必报其仇。乃详为好书遗平原君曰。﹁寡人闻君之高义。原与君为布衣之友。君幸过寡人。寡人原与君为十日之饮。﹂平原君畏秦。且以为然。而入秦见昭王。昭王与平原君饮数日。昭王谓平原君曰。﹁昔周文王得吕尚以为太公。齐桓公得管夷吾以为仲父。今范君亦寡人之叔父也。范君之仇在君之家。原使人归取其头来。不然。吾不出君於关。﹂平原君曰。﹁贵而为交者。为贱也。富而为交者。为贫也。夫魏齐者。胜之友也。在。固不出也。今又不在臣所。﹂昭王乃遗赵王书曰。﹁王之弟在秦。范君之仇魏齐在平原君之家。王使人疾持其头来。不然。吾举兵而伐赵。又不出王之弟於关。﹂赵孝成王乃发卒围平原君家。急。魏齐夜亡出。见赵相虞卿。虞卿度赵王终不可说。乃解其相印。与魏齐亡。间行。念诸侯莫可以急抵者。乃复走大梁。欲因信陵君以走楚。信陵君闻之。畏秦。犹豫未肯见。曰。﹁虞卿何如人也。﹂时侯嬴在旁。曰。﹁人固未易知。知人亦未易也。夫虞卿蹑屩檐簦。一见赵王。赐白璧一双。黄金百镒。再见。拜为上卿。三见。卒受相印。封万户侯。当此之时。天下争知之。夫魏齐穷困过虞卿。虞卿不敢重爵禄之尊。解相印。捐万户侯而间行。急士之穷而归公子。公子曰﹃何如人﹄。人固不易知。知人亦未易也。﹂信陵君大惭。驾如野迎之。魏齐闻信陵君之初难见之。怒而自刭。赵王闻之。卒取其头予秦。秦昭王乃出平原君归赵。\\
\hspace*{2em}昭王四十三年。秦攻韩汾陉。拔之。因城河上广武。\\
\hspace*{2em}後五年。昭王用应侯谋。纵反间卖赵。赵以其故。令马服子代廉颇将。秦大破赵於长平。遂围邯郸。已而与武安君白起有隙。言而杀之。任郑安平。使击赵。郑安平为赵所围。急。以兵二万人降赵。应侯席稿请罪。秦之法。任人而所任不善者。各以其罪罪之。於是应侯罪当收三族。秦昭王恐伤应侯之意。乃下令国中。﹁有敢言郑安平事者。以其罪罪之。﹂而加赐相国应侯食物日益厚。以顺適其意。後二岁。王稽为河东守。与诸侯通。坐法诛。而应侯日益以不怿。\\
\hspace*{2em}昭王临朝叹息。应侯进曰。﹁臣闻﹃主忧臣辱。主辱臣死﹄。今大王中朝而忧。臣敢请其罪。﹂昭王曰。﹁吾闻楚之铁剑利而倡优拙。夫铁剑利则士勇。倡优拙则思虑远。夫以远思虑而御勇士。吾恐楚之图秦也。夫物不素具。不可以应卒。今武安君既死。而郑安平等畔。内无良将而外多敌国。吾是以忧。﹂欲以激励应侯。应侯惧。不知所出。蔡泽闻之。往入秦也。\\
\hspace*{2em}蔡泽者。燕人也。游学干诸侯小大甚众。不遇。而从唐举相。曰。﹁吾闻先生相李兑。曰﹃百日之内持国秉﹄。有之乎。﹂曰。﹁有之。﹂曰。﹁若臣者何如。﹂唐举孰视而笑曰。﹁先生曷鼻。巨肩。魋颜。蹙齃。膝挛。吾闻圣人不相。殆先生乎。﹂蔡泽知唐举戏之。乃曰。﹁富贵吾所自有。吾所不知者寿也。原闻之。﹂唐举曰。﹁先生之寿。从今以往者四十三岁。﹂蔡泽笑谢而去。谓其御者曰。﹁吾持粱刺齿肥。跃马疾驱。怀黄金之印。结紫绶於要。揖让人主之前。食肉富贵。四十三年足矣。﹂去之赵。见逐。之韩。魏。遇夺釜鬲於涂。闻应侯任郑安平。王稽皆负重罪於秦。应侯内惭。蔡泽乃西入秦。\\
\hspace*{2em}将见昭王。使人宣言以感怒应侯曰。﹁燕客蔡泽。天下雄俊弘辩智士也。彼一见秦王。秦王必困君而夺君之位。﹂应侯闻。曰。﹁五帝三代之事。百家之说。吾既知之。众口之辩。吾皆摧之。是恶能困我而夺我位乎。﹂使人召蔡泽。蔡泽入。则揖应。应侯固不快。及见之。又倨。应侯因让之曰。﹁子尝宣言欲代我相秦。宁有之乎。﹂对曰。﹁然。﹂应侯曰。﹁请闻其说。﹂蔡泽曰。﹁吁。君何见之晚也。夫四时之序。成功者去。夫人生百体坚彊。手足便利。耳目聪明而心圣智。岂非士之原与。﹂应侯曰。﹁然。﹂蔡泽曰。﹁质仁秉义。行道施德。得志於天下。天下怀乐敬爱而尊慕之。皆原以为君王。岂不辩智之期与。﹂应侯曰。﹁然。﹂蔡泽复曰。﹁富贵显荣。成理万物。使各得其所。性命寿长。终其天年而不夭伤。天下继其统。守其业。传之无穷。名实纯粹。泽流千里。世世称之而无绝。与天地终始。岂道德之符而圣人所谓吉祥善事者与。﹂应侯曰。﹁然。﹂\\
\hspace*{2em}蔡泽曰。﹁若夫秦之商君。楚之吴起。越之大夫种。其卒然亦可原与。﹂应侯知蔡泽之欲困己以说。复谬曰。﹁何为不可。夫公孙鞅之事孝公也。极身无贰虑。尽公而不顾私。设刀锯以禁奸邪。信赏罚以致治。披腹心。示情素。蒙怨咎。欺旧友。夺魏公子卬。安秦社稷。利百姓。卒为秦禽将破敌。攘地千里。吴起之事悼王也。使私不得害公。谗不得蔽忠。言不取苟合。行不取苟容。不为危易行。行义不辟难。然为霸主强国。不辞祸凶。大夫种之事越王也。主虽困辱。悉忠而不解。主虽绝亡。尽能而弗离。成功而弗矜。贵富而不骄怠。若此三子者。固义之至也。忠之节也。是故君子以义死难。视死如归。生而辱不如死而荣。士固有杀身以成名。虽义之所在。虽死无所恨。何为不可哉。﹂\\
\hspace*{2em}蔡泽曰。﹁主圣臣贤。天下之盛福也。君明臣直。国之福也。父慈子孝。夫信妻贞。家之福也。故比干忠而不能存殷。子胥智而不能完吴。申生孝而晋国乱。是皆有忠臣孝子。而国家灭乱者。何也。无明君贤父以听之。故天下以其君父为僇辱而怜其臣子。今商君。吴起。大夫种之为人臣。是也。其君。非也。故世称三子致功而不见德。岂慕不遇世死乎。夫待死而後可以立忠成名。是微子不足仁。孔子不足圣。管仲不足大也。夫人之立功。岂不期於成全邪。身与名俱全者。上也。名可法而身死者。其次也。名在僇辱而身全者。下也。﹂於是应侯称善。\\
\hspace*{2em}蔡泽少得间。因曰。﹁夫商君。吴起。大夫种。其为人臣尽忠致功则可原矣。闳夭事文王。周公辅成王也。岂不亦忠圣乎。以君臣论之。商君。吴起。大夫种其可原孰与闳夭。周公哉。﹂应侯曰。﹁商君。吴起。大夫种弗若也。﹂蔡泽曰。﹁然则君之主慈仁任忠。惇厚旧故。其贤智与有道之士为胶漆。义不倍功臣。孰与秦孝公。楚悼王。越王乎。﹂应侯曰。﹁未知何如也。﹂蔡泽曰。﹁今主亲忠臣。不过秦孝公。楚悼王。越王。君之设智。能为主安危修政。治乱彊兵。批患折难。广地殖穀。富国足家。彊主。尊社稷。显宗庙。天下莫敢欺犯其主。主之威盖震海内。功彰万里之外。声名光辉传於千世。君孰与商君。吴起。大夫种。﹂应侯曰。﹁不若。﹂蔡泽曰。﹁今主之亲忠臣不忘旧故不若孝公。悼王。句践。而君之功绩爱信亲幸又不若商君。吴起。大夫种。然而君之禄位贵盛。私家之富过於三子。而身不退者。恐患之甚於三子。窃为君危之。语曰﹃日中则移。月满则亏﹄。物盛则衰。天地之常数也。进退盈缩。与时变化。圣人之常道也。故﹃国有道则仕。国无道则隐﹄。圣人曰﹃飞龙在天。利见大人﹄。﹃不义而富且贵。於我如浮云﹄。今君之怨已雠而德已报。意欲至矣。而无变计。窃为君不取也。且夫翠。鹄。犀。象。其处势非不远死也。而所以死者。惑於饵也。苏秦。智伯之智。非不足以辟辱远死也。而所以死者。惑於贪利不止也。是以圣人制礼节欲。取於民有度。使之以时。用之有止。故志不溢。行不骄。常与道俱而不失。故天下承而不绝。昔者齐桓公九合诸侯。一匡天下。至於葵丘之会。有骄矜之志。畔者九国。吴王夫差兵无敌於天下。勇彊以轻诸侯。陵齐晋。故遂以杀身亡国。夏育。太史噭叱呼骇三军。然而身死於庸夫。此皆乘至盛而不返道理。不居卑退处俭约之患也。夫商君为秦孝公明法令。禁奸本。尊爵必赏。有罪必罚。平权衡。正度量。调轻重。决裂阡陌。以静生民之业而一其俗。劝民耕农利土。一室无二事。力田稸积。习战陈之事。是以兵动而地广。兵休而国富。故秦无敌於天下。立威诸侯。成秦国之业。功已成矣。而遂以车裂。楚地方数千里。持戟百万。白起率数万之师以与楚战。一战举鄢郢以烧夷陵。再战南并蜀汉。又越韩。魏而攻彊赵。北阬马服。诛屠四十馀万之众。尽之于长平之下。流血成川。沸声若雷。遂入围邯郸。使秦有帝业。楚。赵天下之彊国而秦之仇敌也。自是之後。楚。赵皆慑伏不敢攻秦者。白起之势也。身所服者七十馀城。功已成矣。而遂赐剑死於杜邮。吴起为楚悼王立法。卑减大臣之威重。罢无能。废无用。损不急之官。塞私门之请。一楚国之俗。禁游客之民。精耕战之士。南收杨越。北并陈。蔡。破横散从。使驰说之士无所开其口。禁朋党以励百姓。定楚国之政。兵震天下。威服诸侯。功已成矣。而卒枝解。大夫种为越王深谋远计。免会稽之危。以亡为存。因辱为荣。垦草入邑。辟地殖穀。率四方之士。专上下之力。辅句践之贤。报夫差之雠。卒擒劲吴。令越成霸。功已彰而信矣。句践终负而杀之。此四子者。功成不去。祸至於此。此所谓信而不能诎。往而不能返者也。范蠡知之。超然辟世。长为陶硃公。君独不观夫博者乎。或欲大投。或欲分功。此皆君之所明知也。今君相秦。计不下席。谋不出廊庙。坐制诸侯。利施三川。以实宜阳。决羊肠之险。塞太行之道。又斩范。中行之涂。六国不得合从。栈道千里。通於蜀汉。使天下皆畏秦。秦之欲得矣。君之功极矣。此亦秦之分功之时也。如是而不退。则商君。白公。吴起。大夫种是也。吾闻之。﹃鉴於水者见面之容。鉴於人者知吉与凶﹄。书曰﹃成功之下。不可久处﹄。四子之祸。君何居焉。君何不以此时归相印。让贤者而授之。退而岩居川观。必有伯夷之廉。长为应侯。世世称孤。而有许由。延陵季子之让。乔松之寿。孰与以祸终哉。即君何居焉。忍不能自离。疑不能自决。必有四子之祸矣。易曰﹃亢龙有悔﹄。此言上而不能下。信而不能诎。往而不能自返者也。原君孰计之。﹂应侯曰。善。吾闻﹃欲而不知。失其所以欲。有而不知足。失其所以有﹄。先生幸教。雎敬受命。於是乃延入坐。为上客。\\
\hspace*{2em}後数日。入朝。言於秦昭王曰。﹁客新有从山东来者曰蔡泽。其人辩士。明於三王之事。五伯之业。世俗之变。足以寄秦国之政。臣之见人甚众。莫及。臣不如也。臣敢以闻。﹂秦昭王召见。与语。大说之。拜为客卿。应侯因谢病请归相印。昭王彊起应侯。应侯遂称病笃。范雎免相。昭王新说蔡泽计画。遂拜为秦相。东收周室。\\
\hspace*{2em}蔡泽相秦数月。人或恶之。惧诛。乃谢病归相印。号为纲成君。居秦十馀年。事昭王。孝文王。庄襄王。卒事始皇帝。为秦使於燕。三年而燕使太子丹入质於秦。\\
\hspace*{2em}太史公曰。韩子称﹁长袖善舞。多钱善贾﹂。信哉是言也。范雎。蔡泽世所谓一切辩士。然游说诸侯至白首无所遇者。非计策之拙。所为说力少也。及二人羁旅入秦。继踵取卿相。垂功於天下者。固彊弱之势异也。然士亦有偶合。贤者多如此二子。不得尽意。岂可胜道哉。然二子不困戹。恶能激乎。\\

% ===== 索隐述赞（司马贞《史记索隐》）· 统一样式 =====
% 题头「索隐述赞曰」五字并入反白框；正文每四句一列，定宽盒 \hfil 居列中；
% \fontsize{24}{41} 列距对齐全书界行（121+41k）；\clearpage 使述赞独占一页。
\clearpage
\shuzanhead{索隐述赞曰}\linebreak
{\shuzanbody\fontsize{24}{41}\selectfont
\hbox to \hsize{\hfil 应侯始困。讬载而西。说行计立。贵平宠稽。\hfil}\linebreak
\hbox to \hsize{\hfil 倚秦市赵。卒报魏齐。纲成辩智。范雎招携。\hfil}\linebreak
\hbox to \hsize{\hfil 势利倾夺。一言成蹊。\hfil}%
\par}
\end{document}
